% Source: Weinberg GR, physical PDF pp. 65--67
%         (printed pp. 39--41).
% Coverage: Section 2.6, Currents and Densities; equations
%           (2.6.1)--(2.6.8).
% Figures/tables/footnotes: none.
% Status: source-reviewed and compile-clean.
% Uncertainties: none.

\subsection{Currents and Densities}\label{sec:2.6}

Suppose that we have a system of particles with positions
\(\mathbf{x}_n(t)\) and charges \(e_n\). The current and charge densities are
usually defined by
\begin{equation}
  \mathbf J(\mathbf x,t)
  =\sum_n e_n\delta^3\bigl(\mathbf x-\mathbf x_n(t)\bigr)
  \frac{\dd\mathbf x_n(t)}{\dd t},
  \tag{2.6.1}\label{eq:2.6.1}
\end{equation}
\begin{equation}
  \rho(\mathbf x,t)
  =\sum_n e_n\delta^3\bigl(\mathbf x-\mathbf x_n(t)\bigr).
  \tag{2.6.2}\label{eq:2.6.2}
\end{equation}
Here \(\delta^3\) is the Dirac delta function, defined by the statement that
for any smooth function \(f(\mathbf x)\),
\[
  \int \dd^3x\,f(\mathbf x)\delta^3(\mathbf x-\mathbf y)=f(\mathbf y).
\]
We can unite \(\mathbf J\) and \(\rho\) into a four-vector \(J^\mu\) by
setting
\begin{equation}
  J^0=\rho,
  \tag{2.6.3}\label{eq:2.6.3}
\end{equation}
that is,
\begin{equation}
  J^\mu(x)
  =\sum_n e_n\delta^3\bigl(\mathbf x-\mathbf x_n(t)\bigr)
  \frac{\dd x_n^\mu(t)}{\dd t}.
  \tag{2.6.4}\label{eq:2.6.4}
\end{equation}
To show that this is a four-vector, define \(x_n^0(t)=t\), and write
Eq.~\eqref{eq:2.6.4} as
\[
  J^\mu(x)=\int \dd t'\sum_n e_n
  \delta^4\bigl(x-x_n(t')\bigr)\frac{\dd x_n^\mu(t')}{\dd t'} .
\]
The differentials \(\dd t'\) cancel, and hence can be replaced with an invariant
\(\dd\tau\):
\begin{equation}
  J^\mu(x)=\int \dd\tau\sum_n e_n
  \delta^4\bigl(x-x_n(\tau)\bigr)\frac{\dd x_n^\mu(\tau)}{\dd\tau}.
  \tag{2.6.5}\label{eq:2.6.5}
\end{equation}
But \(\delta^4(x-x_n(\tau))\) is a scalar (because
\(\det\Lambda=1\)) and \(\dd x_n^\mu\) is a four-vector, so \(J^\mu\) is a
four-vector.

We also note that
\[
\begin{aligned}
  \boldsymbol{\nabla}\cdot\mathbf J(\mathbf x,t)
  &=\sum_n e_n
    \partial_i
    \delta^3\bigl(\mathbf x-\mathbf x_n(t)\bigr)
    \frac{\dd x_n^i(t)}{\dd t}\\
  &=-\sum_n e_n
    % NOTATION EXCEPTION: the differentiated particle coordinate must remain
    % explicit to distinguish it from the field coordinate in the prior line.
    \frac{\partial}{\partial x_n^i}
    \delta^3\bigl(\mathbf x-\mathbf x_n(t)\bigr)
    \frac{\dd x_n^i(t)}{\dd t}\\
  &=-\sum_n e_n\partial_t
    \delta^3\bigl(\mathbf x-\mathbf x_n(t)\bigr)\\
  &=-\partial_t\rho(\mathbf x,t),
\end{aligned}
\]
or, in four-dimensional language,
\begin{equation}
  \partial_\mu J^\mu(x)=0.
  \tag{2.6.6}\label{eq:2.6.6}
\end{equation}
The Lorentz invariance of this statement is evident.

Whenever any current \(J^\mu(x)\) satisfies the invariant conservation law
\eqref{eq:2.6.6}, we can form a total charge
\begin{equation}
  Q=\int \dd^3x\,J^0(x).
  \tag{2.6.7}\label{eq:2.6.7}
\end{equation}
This quantity is time-independent, because Eq.~\eqref{eq:2.6.6} and Gauss's
theorem give
\[
  \frac{\dd Q}{\dd t}
  =\int \dd^3x\,\partial_0J^0(x)
  =-\int \dd^3x\,\boldsymbol{\nabla}\cdot\mathbf J(x)=0 .
\]
If \(J^\mu(x)\) is a four-vector, then \(Q\) is not only constant but a
scalar. To see this, write \(Q\) as
\begin{equation}
  Q=\int \dd^4x\,J^\mu(x)\,
  \partial_\mu\theta(n_\nu x^\nu),
  \tag{2.6.8}\label{eq:2.6.8}
\end{equation}
where \(\theta\) is the step function
\[
  \theta(s)=
  \begin{cases}
    1,&s>0,\\
    0,&s<0,
  \end{cases}
\]
and \(n_\mu\) is defined by
\[
  n_0=1,\qquad n_1=n_2=n_3=0 .
\]
The effect of a Lorentz transformation on \(Q\) is then evidently simply to
change \(n_\mu\):
\[
  Q'=\int \dd^4x\,J^\mu(x)\,
  \partial_\mu\theta(n'_\nu x^\nu),
  \qquad
  n'_\mu=\Lambda_\mu{}^\nu n_\nu ,
\]
and using Eq.~\eqref{eq:2.6.6}, the change in \(Q\) is then
\[
  Q'-Q
  =\int \dd^4x\,J^\mu(x)\partial_\mu
  \left[
    \theta(n'_\nu x^\nu)-\theta(n_\nu x^\nu)
  \right].
\]
The current \(J^\mu(x)\) can be presumed to vanish if
\(\lvert\mathbf x\rvert\to\infty\) with \(t\) fixed, whereas the function
\(\theta(n'_\nu x^\nu)-\theta(n_\nu x^\nu)\) vanishes if
\(\lvert t\rvert\to\infty\) with \(\mathbf x\) fixed. Hence we can apply the
four-dimensional Gauss theorem, and find \(Q'-Q=0\); that is, \(Q\) is a
scalar.

(For the current density \(J^\mu\) defined by Eq.~\eqref{eq:2.6.4}, the
charge \eqref{eq:2.6.7} is
\[
  Q=\sum_n e_n ,
\]
which of course is a constant scalar; however, in dealing with the charge and
current distributions of extended particles it is important to realize that
Eq.~\eqref{eq:2.6.7} defines a time-independent scalar for any conserved
four-vector \(J^\mu\).)
